\section{问题重述}

\paperexample{实际论文在此说明研究对象、数据或业务背景，以及题面要求的输入、输出、评价目标和限制。不要在背景段落中引入题面没有授权的结论。}

对于多子问题任务，建议先建立一张问题契约表：每个子问题的输入、输出、决策或响应变量、核心目标、硬约束、评价指标、证据位置和负责人均应明确。题面和公开资料的来源应在正文中按顺序编码制引用；本模板的官方格式来源见文献 \cite{cmathc2025format,cmathc2025notice}。

\begin{table}[H]
  \centering
  \caption{问题契约示例}
  \label{tab:problem-contract}
  \begin{tabular}{p{0.16\textwidth}p{0.22\textwidth}p{0.22\textwidth}p{0.28\textwidth}}
    \toprule
    子问题 & 输入 & 输出 & 核心目标与约束 \\
    \midrule
    子问题一 & 数据/参数 & 统计或决策结果 & 指标、边界、证据 \\
    子问题二 & 上游结果 & 模型/策略 & 资源、时间、可行性 \\
    \bottomrule
  \end{tabular}
\end{table}
